% Part: first-order-logic
% Chapter: syntax-and-semantics
% Section: formation-sequences

\documentclass[../../../include/open-logic-section]{subfiles}

\begin{document}

\olfileid{pl}{syn}{fseq}

\olsection{دنباله‌های ساخت}

تعریف استقراییِ !!{formula}s و فن مکملِ اثبات ویژگی‌های !!{formula}s با
استقرا، رویکردی ظریف و کارآمد است. بااین‌همه، بررسی رویکردی از پایین به بالا
و گام‌به‌گام برای ساخت !!{formula}s نیز می‌تواند سودمند باشد؛ در اینجا این
کار را با استفاده از مفهوم یک \emph{دنبالهٔ ساخت} انجام می‌دهیم.

\begin{defn}[دنباله‌های ساخت برای فرمول‌ها]
\ollabel{defn:fseq-frm}
دنبالهٔ متناهیِ $\tuple{!A_0,\dotsc,!A_n}$ از رشته‌های نمادهای
زبان~$\Lang L_0$ یک \emph{دنبالهٔ ساخت} برای $!A$ است اگر
$!A \ident !A_n$ و برای هر $i \leq n$، یا $!A_i$ فرمولی اتمی باشد یا
$j,k < i$ وجود داشته باشند، به‌گونه‌ای که یکی از موارد زیر برقرار باشد:
\begin{enumerate}
    \tagitem{prvNot}{$!A_i \ident \lnot !A_j$.}{}%
    \tagitem{prvAnd}{$!A_i \ident (!A_j \land !A_k)$.}{}%
    \tagitem{prvOr}{$!A_i \ident (!A_j \lor !A_k)$.}{}%
    \tagitem{prvIf}{$!A_i \ident (!A_j \lif !A_k)$.}{}%
    \tagitem{prvIff}{$!A_i \ident (!A_j \liff !A_k)$.}{}%
\end{enumerate}
\end{defn}

\begin{ex}
\[
\tuple{
    \Obj p_0,
    \Obj p_1,
    (\Obj p_1 \land \Obj p_0),
    \lnot (\Obj p_1 \land \Obj p_0)
}
\]
یک دنبالهٔ ساخت برای
$\lnot (\Obj p_1 \land \Obj p_0)$ است؛ دنبالهٔ زیر نیز چنین است:
\[
\tuple{
    \Obj p_0,
    \Obj p_1,
    \Obj p_0,
    (\Obj p_1 \land \Obj p_0),
    (\Obj p_0 \lif \Obj p_1),
    \lnot (\Obj p_1 \land \Obj p_0)
}.
\]
%
چنان‌که از مثال دوم پیداست، دنباله‌های ساخت ممکن است حاوی «اجزای زائد»
باشند: فرمول‌هایی که تکراری‌اند یا در ساخت نقشی ندارند.
\end{ex}

\begin{prop}
\ollabel{prop:formed}
هر !!{formula}~$!A$ در~$\Frm[L_0]$ دارای یک دنبالهٔ ساخت است.
\end{prop}

\begin{proof}
فرض کنید $!A$ اتمی باشد. آنگاه دنبالهٔ~$\tuple{!A}$ یک دنبالهٔ ساخت
برای~$!A$ است.
%
اکنون فرض کنید $!B$ و~$!C$ به‌ترتیب دارای دنباله‌های ساخت
$\tuple{!B_0,\dotsc,!B_n}$ و $\tuple{!C_0,\dotsc,!C_m}$ باشند.
%
\begin{enumerate}
  \tagitem{prvNot}{اگر $!A \ident \lnot !B$،
      آنگاه $\tuple{!B_0,\dotsc,!B_n,\lnot !B_n}$
      یک دنبالهٔ ساخت برای~$!A$ است.}{}
  \tagitem{prvAnd}{اگر $!A \ident (!B \land !C)$،
      آنگاه $\tuple{!B_0,\dotsc,!B_n,!C_0,\dotsc,!C_m,(!B_n \land !C_m)}$
      یک دنبالهٔ ساخت برای~$!A$ است.}{}
  \tagitem{prvOr}{اگر $!A \ident (!B \lor !C)$،
      آنگاه $\tuple{!B_0,\dotsc,!B_n,!C_0,\dotsc,!C_m,(!B_n \lor !C_m)}$
      یک دنبالهٔ ساخت برای~$!A$ است.}{}
  \tagitem{prvIf}{اگر $!A \ident (!B \lif !C)$،
      آنگاه $\tuple{!B_0,\dotsc,!B_n,!C_0,\dotsc,!C_m,(!B_n \lif !C_m)}$
      یک دنبالهٔ ساخت برای~$!A$ است.}{}
  \tagitem{prvIff}{اگر $!A \ident (!B \liff !C)$،
      آنگاه $\tuple{!B_0,\dotsc,!B_n,!C_0,\dotsc,!C_m,(!B_n \liff !C_m)}$
      یک دنبالهٔ ساخت برای~$!A$ است.}{}
\end{enumerate}
بنا بر اصل استقرا بر !!{formula}s،
هر !!{formula} دارای یک دنبالهٔ ساخت است.
\end{proof}

عکس این گزاره را نیز می‌توانیم اثبات کنیم. این امر مهم است، زیرا نشان می‌دهد
دو شیوهٔ ما برای تعریف فرمول‌ها هم‌ارزند: نتایج یکسانی به دست می‌دهند. همچنین
بدین معناست که می‌توانیم قضیه‌های مربوط به فرمول‌ها را با استقرای معمولی بر
طول دنباله‌های ساخت اثبات کنیم.

\begin{lem}
\ollabel{lem:fseq-init}
فرض کنید $\tuple{!A_0,\dotsc,!A_n}$ دنبالهٔ ساختی برای~$!A_n$ باشد و
$k \leq n$. آنگاه $\tuple{!A_0,\dotsc,!A_k}$ دنبالهٔ ساختی
برای~$!A_k$ است.
\end{lem}

\begin{proof}
تمرین.
\end{proof}

\begin{thm}
\ollabel{thm:fseq-frm-equiv}
$\Frm[L_0]$ مجموعهٔ همهٔ رشته‌های نمادهای زبان~$\Lang L_0$ است که
دارای دنبالهٔ ساخت‌اند.
\end{thm}

\begin{proof}
بگذارید $F$ مجموعهٔ همهٔ رشته‌های نمادهای زبان~$\Lang L_0$ باشد که
دارای دنبالهٔ ساخت‌اند. در \olref[pl][syn][fseq]{prop:formed} دیدیم که
$\Frm[L_0] \subseteq F$؛ پس اکنون شمول عکس را اثبات می‌کنیم.

فرض کنید $!A$ دارای دنبالهٔ ساخت $\tuple{!A_0,\dotsc,!A_n}$ باشد.
اثبات می‌کنیم که $!A \in \Frm[L_0]$؛ این کار را با استقرای قوی
بر~$n$ انجام می‌دهیم.
فرض استقرای ما این است که هر رشته از نمادها که دنبالهٔ ساختی با شاخص پایانی
$m < n$ دارد، در~$\Frm[L_0]$ است. بنا بر تعریف دنبالهٔ ساخت، یا $!A_n$
اتمی است یا باید $j,k < n$ وجود داشته باشند، به‌گونه‌ای که یکی از موارد زیر
برقرار باشد:
\begin{enumerate}
    \tagitem{prvNot}{$!A_n \ident \lnot !A_j$.}{}%
    \tagitem{prvAnd}{$!A_n \ident (!A_j \land !A_k)$.}{}%
    \tagitem{prvOr}{$!A_n \ident (!A_j \lor !A_k)$.}{}%
    \tagitem{prvIf}{$!A_n \ident (!A_j \lif !A_k)$.}{}%
    \tagitem{prvIff}{$!A_n \ident (!A_j \liff !A_k)$.}{}%
\end{enumerate}
اکنون به تفکیک حالت‌ها استدلال می‌کنیم. اگر $!A_n$ اتمی باشد، آنگاه
$!A_n \in \Frm[L_0]$. حال فرض کنید $!A \ident
(!A_j \land !A_k)$. بنا بر \olref[pl][syn][fseq]{lem:fseq-init}،
$\tuple{!A_0,\dotsc,!A_j}$ و $\tuple{!A_0,\dotsc,!A_k}$ به‌ترتیب
دنباله‌های ساختی برای $!A_j$ و~$!A_k$ هستند. چون هر یک قطعهٔ آغازینِ
سره‌ای از دنبالهٔ ساختِ~$!A$ است، شاخص پایانیِ هر دو کوچک‌تر از~$n$ است.
بنابراین، بنا بر فرض استقرا، $!A_j$ و~$!A_k$ در~$\Frm[L_0]$ هستند و
ازاین‌رو، بنا بر تعریف !!a{formula}،
$(!A_j \land !A_k)$ نیز چنین است. حالت‌های دیگر با استدلالی موازی
به دست می‌آیند.
\end{proof}

\end{document}
